\section{Methodik}
Diese Arbeit folgt einem mehrstufigen methodischen Vorgehen, das sowohl theoretische als auch praktische Forschungselemente kombiniert. Ziel ist es, die Einsatzmöglichkeiten von \ac{KI} in \ac{ERP}-Systemen zu analysieren und darauf aufbauend einen \ac{KI}-Agenten in \ac{D365FnSCM} prototypisch zu entwickeln.

Zunächst wird durch eine klassische Literaturrecherche der theoretische Rahmen der Arbeit erarbeitet. Darauf aufbauend erfolgt eine systematische Literaturrecherche nach dem PRISMA-Vorgehen zum Einsatz von \ac{KI} in \ac{ERP}-Systemen.

Auf Basis der gewonnenen Erkenntnisse wird anschließend das Design Science Research-Paradigma (DSR) nach Österle angewendet, um ein IT-Artefakt in Form eines \ac{KI}-Agenten zu konzipieren, zu implementieren und hinsichtlich seiner Effizienzpotenziale zu evaluieren.

\subsection{Literaturrecherche}
Die Literaturrecherche dieser Arbeit erfolgt in zwei Schritten, bestehend aus einer explorativen und einer systematischen Literaturrecherche.

Zunächst wird eine explorative Literaturrecherche durchgeführt, um die theoretischen Grundlagen der in dieser Arbeit verwendeten zentralen Konzepte zu erarbeiten. Diese Recherchephase diente der Klärung grundlegender Begriffe und Modelle in den Bereichen Enterprise-Resource-Planning-Systeme und Künstliche Intelligenz sowie deren jeweilige Teilbereiche und Ausprägungen. Ziel war es, ein konsistentes Begriffsverständnis und einen stabilen theoretischen Rahmen zu schaffen, auf dem die restliche Arbeit aufbaut. Eine solche explorative Vorgehensweise eignet sich in diesem Fall besonders, da mehrere komplexe Konzepte aus unterschiedlichen Disziplinen zusammengeführt werden. Snyder weist darauf hin, dass systematische Literaturrecherchen vor allem für klar abgegrenzte Fragestellungen geeignet sind, während bei breiten oder heterogenen Themenfeldern zunächst eine vorgelagerte Orientierungs- und Strukturierungsphase erforderlich ist, um Forschungsgegenstand eindeutig zu definieren \cite[\pagef 334ff]{snyder_literature_2019}.

Aufbauend auf den Ergebnissen der explorativen Literaturrecherche wird eine systematische Literaturrecherche durchgeführt, die sich gezielt auf den Einsatz von Künstlicher Intelligenz in ERP-Systemen konzentriert. Aufgrund der nun klar abgegrenzten Fragestellung eignet sich hierfür ein systematisches Vorgehen.

Für die systematische Literaturrecherche wird das \ac{PRISMA}-Framework herangezogen. PRISMA stellt einen Standard dar, der eine klar strukturierte und nachvollziehbare Methodik zur Identifikation, Auswahl und Dokumentation relevanter wissenschaftlicher Literatur vorgibt. Durch die transparente Definition von Suchstrategien sowie expliziter Ein- und Ausschlusskriterien wird eine fundierte, reproduzierbare und methodisch belastbare Grundlage für die weitere wissenschaftliche Analyse geschaffen. \cite[\pagef n71]{page_prisma_2021}
\subsection{Design Science Research}
Aufbauend auf den Erkenntnissen aus der Literaturrecherche wird zur Entwicklung und Untersuchen eines \ac{KI}-Agentendas Vorgehensmodell des \ac{DSR} angewendet, 
da der Schwerpunkt auf der Entwicklung und prototypischen Implementierung eines \ac{KI}-Agenten in einem \ac{ERP}-System liegt. \ac{DSR} ist ein problemorientierter Forschungsansatz, der darauf ausgelegt ist, IT-Artefakte zu entwickeln und zu bewerten, um konkrete Probleme im Bereich der Informationssysteme zu adressieren. \ac{DSR} ist insbesondere dann geeignet, wenn informationssystembasierte Anwendungen entwickelt werden, um organisatorische Prozesse zu unterstützen oder zu verbessern. \cite[\pagef 1ff]{hevner_design_2010}
Der in dieser Arbeit entwickelte \ac{KI}-Agent für \ac{D365FnSCM} stellt genau ein solches IT-Artefakt dar, weshalb \ac{DSR} sich somit ideal für die methodische Vorgehensweise eignet.

